\chapter{Is Ray Tune the right base, and what we should build on it}
\label{sec:architecture}

The question that started this document can now be answered with the
evidence laid out. Yes, Ray Tune is the right base, provided we are clear
about which layer of the problem it solves. Tune should run our trials; it
should not choose them. The search intelligence, which is where the field
has moved fastest and where Tune has moved least, belongs in a layer we
own. This chapter turns that sentence into a design, and the design is
small: the whole proposal is two interfaces, two adapters, and a database.

\section{The case for Tune underneath}

Three facts carry the decision. First, we already operate Ray: RLlib
training runs on the same cluster, speaks the same actor model, and
produces the checkpoints that population-based methods need to copy and
resume. A second execution stack would be pure overhead. Second, Tune's
execution machinery is the strongest available in open source for our
shape of work: multi-node trial scheduling, fault tolerance,
checkpoint-aware schedulers, and the only mainstream native distributed
PBT and PB2 (Chapter~\ref{sec:libraries}). The population family of
Chapter~\ref{sec:population} is the part of our roadmap with the most
demanding executor requirements, exploit steps are checkpoint surgery
performed live on a running population, and Tune is the one framework
that has already productionized them. Third, the composition we propose
is the one the ecosystem itself converges on: Optuna and Ray document
each other as the intended pairing, a 2025 industrial optimizer built
its stack exactly this way, and Google's published account of its own
tuning service benchmarks against algorithms accessed through Ray Tune
\citep{song2024vizier}.

The risks deserve equal daylight, and Chapter~\ref{sec:libraries}
documented them with dates. Tune's own search layer has been coasting:
across the five Ray releases from late 2025 through August 2026, every
Tune change was a fix or a documentation edit, several of its remaining
searcher integrations wrap upstreams that are frozen or formally
unmaintained, and core Tune still has no native concept of a Pareto
front, which after Chapter~\ref{sec:moo} we regard as a first-class
requirement rather than a nicety. None of this argues against Tune as an
executor; all of it argues against depending on Tune for algorithms.
There is also a real overhead concern at the small end: for the cheap
trials of our Hamiltonian-network work
(Section~\ref{sec:costmodel}), cluster machinery can cost more than the
training it schedules.

\section{The design}

The consequence is a three-layer design, drawn in
Figure~\ref{fig:architecture}. At the top, a search layer that is the
actual product: searchers (ask-and-tell proposers: GP with qLogEI and
qLogNEHVI on BoTorch, TPE via Optuna, CMA-ES, the mixed-space
trust-region searcher that BG-PBT needs) and schedulers (ASHA and the
median rule, their multi-objective variants with veto gates, the
population controllers), all written against two small interfaces of our
own. The split between the two interfaces is exactly the split drawn in
Figure~\ref{fig:loop} at the start of the document: searchers answer
``which configuration next,'' schedulers answer ``how much budget, and
does this trial live''; every method surveyed decomposes into the two
roles, and keeping the roles separate is what lets ASHA schedule trials
proposed by any searcher.

In the middle, thin adapters: one maps our interfaces onto Ray Tune's
searcher and trial-scheduler protocols and is the default; one runs
trials in local processes for laptop-scale studies where Ray does not
earn its overhead; the design leaves room for a queue-based backend
later, following the backend abstraction that Syne Tune demonstrated
\citep{salinas2022synetune}. At the bottom, alongside the executors, a
run store that records every trial of every study: configuration, seeds,
fidelity-indexed curves, cost, diagnostics, and lineage (which weights
were copied from whom, for population methods). The store is what turns
the transfer methods of Chapter~\ref{sec:transfer} from papers into
features, warm starts are a query against it, and $\pi$BO priors and
fitted cost models are trained from it, and it is also where a study's
evidence lives when we later ask whether the tuner is earning its keep.

To make the layering concrete, walk one trial of the running example
through it. The GP searcher (top layer) proposes a configuration by
maximizing qLogEI over the mixed space; the MO-ASHA scheduler assigns it
a rung-1 budget of two million steps. The Ray adapter translates both
decisions into a Tune trial with a checkpoint interval; RLlib trains it
on the cluster and streams evaluation returns back. At the rung
boundary the scheduler ranks the rung (non-dominated sorting if the
study is multi-objective), checks the trial's veto fields, and either
promotes it, at which point the adapter resumes it from its checkpoint,
or stops it. Every intermediate return, the cost, the seed, and the
promotion decision land in the run store. Nothing in the top layer knows
Ray exists; nothing below the adapters knows what an acquisition
function is. That ignorance is the design.

Owning the two interfaces, rather than coding directly against Tune's, is
what keeps the Tune risks bounded. If Tune stagnates further or Ray's
APIs shift again, the adapters change and the product does not. It is
also what lets the same searcher serve both executors, so a Hamiltonian
study on a workstation and a PPO population on the cluster run the same
code paths above the adapter line.

Three features must be first-class in those interfaces because
retrofitting them is painful, and each traces to a specific finding
earlier in the document. Multiple objectives and constraint vetoes: a
scheduler must be able to rank partial results by non-dominated sorting
and to refuse promotion to a trial that failed a correctness diagnostic,
whatever its score; this is the promotion machinery of
Chapter~\ref{sec:moo} and the veto gates that
Figure~\ref{fig:sampler-loops} showed the sampler workload demanding.
Seeds: a trial is a configuration times a seed, aggregation across seeds
is the tuner's job, and tuning seeds are disjoint from test seeds by
construction, the protocol whose absence corrupted the RL results
surveyed in Chapter~\ref{sec:workloads}
\citep{eimer2023hyperparameters}. Schedules: the outcome of a population
study is a trajectory with lineage, storable and replayable, the bold
path of Figure~\ref{fig:pbt} as a first-class object (Tune's own
PBT-replay scheduler shows the shape).

\begin{figure}[t]
\centering
\begin{tikzpicture}[
  box/.style={draw, rounded corners, align=center, inner sep=5pt, font=\small},
  lab/.style={font=\small\itshape, anchor=west},
  node distance=4mm and 5mm]
\node[box, minimum width=136mm, minimum height=17mm] (search) at (0,0)
  {searchers: Sobol, TPE (Optuna), GP + qLogEI / qLogNEHVI (BoTorch), CMA-ES, mixed-space trust region\\
   schedulers: ASHA, median rule, MO-ASHA + veto gates, PBT / PB2 / BG-PBT controllers, PriorBand};
\node[lab] at (-6.8,1.35) {search layer (ours)};
\node[box, minimum width=64mm, minimum height=9mm, below=of search.south west, anchor=north west, xshift=0mm] (ray)
  {Ray Tune adapter (cluster, RLlib, populations)};
\node[box, minimum width=64mm, minimum height=9mm, below=of search.south east, anchor=north east] (local)
  {local adapter (workstation-scale studies)};
\node[box, minimum width=136mm, minimum height=9mm] (store) at ($(ray.south east)!0.5!(local.south west)$) [yshift=-9mm]
  {run store: configurations, seeds, curves, costs, diagnostics, lineage; feeds warm starts and transfer};
\draw[-{Latex}] (search.south -| ray) -- (ray.north);
\draw[-{Latex}] (search.south -| local) -- (local.north);
\draw[-{Latex}] (ray.south) -- (ray.south |- store.north);
\draw[-{Latex}] (local.south) -- (local.south |- store.north);
\end{tikzpicture}
\caption{The proposed layering. Algorithms live above the adapter line
and survive any change of executor; Ray Tune supplies distributed
execution and checkpoint surgery; a shared run store accumulates the
evidence that priors, warm starts, and transfer methods consume. The
two interfaces at the adapter line are the searcher/scheduler split of
Figure~\ref{fig:loop}, made contractual.}
\label{fig:architecture}
\end{figure}

\section{Alternatives we considered}

Adopting Syne Tune outright would buy the broadest method zoo in one
package, and its simulator backend is a genuinely good idea we intend to
copy; against that stand a small maintainer base, no Ray backend, and a
method set that still lacks the BG-PBT-class population line we care
about.

Building on Optuna alone would give us the best samplers and
nothing to run them with: its documented multi-node story is a shared
relational database, and there is no population machinery. NePS has the
freshest prior-and-fidelity algorithms and we will use it as a reference
implementation and correctness oracle, but its filesystem coordination
and small community rule it out as a substrate.

Katib solves execution
for Kubernetes shops; we are a Ray shop. Writing our own executor from
scratch is the option nobody should pick while Tune exists and is
maintained. The remaining risk in our choice is Ray itself drifting in a
direction that breaks Tune; the adapter boundary is the insurance, and
the annual cost of that insurance is two small interfaces.
